% !TeX program = lualatex
% !TeX encoding = utf8
% !TeX spellcheck = uk_UA
% !TeX root = ../main.tex

%=========================================================
\chapter{Електричні заряди та їх взаємодія}\label{\currfilebase}\hypertarget{Charges}{}
\graphicspath{{\currfilebase/Pictures}}
%=========================================================
\epigraph{\Alexander  ...взаємне тяжіння електричної рідини, іменованої позитивною, і електричної рідини, іменованої зазвичай негативною,
	полягає у зворотному відношенні квадратів відстаней...
}{Шарль Огюстен Кулон}


%% --------------------------------------------------------
\section{Сили (взаємодії) в природі}\hypertarget{fundamental-interactions}{}
%% --------------------------------------------------------

Фізичні явища в природі зумовлені взаємодією між тілами. Наприклад, рух Землі навколо Сонця пояснюється гравітаційною взаємодією, яка визначає орбіту
планети. Магнітні сили притягують залізні предмети, а молекулярні сили дозволяють воді підійматися в капілярі. $\beta$-Радіоактивний розпад
ядер є результатом слабкої взаємодії, тоді як притягання між нуклонами в ядрі відбувається завдяки сильній ядерній взаємодії. Попри
різноманіття
цих явищ, усі взаємодії в природі можна звести до чотирьох фундаментальних типів: гравітаційної, електромагнітної, сильної (ядерної) та слабкої
взаємодій.


Гравітаційна та електромагнітна взаємодії досліджуються в рамках класичної фізики, хоча для електромагнітних сил, на відміну від гравітаційних,
розроблена квантова теорія, відома як квантова електродинаміка. Водночас сильна та слабка взаємодії описуються виключно в
рамках квантової теорії. Ці чотири типи взаємодій охоплюють весь спектр фізичних процесів --- від подій у Всесвіті до явищ на рівні елементарних
частинок та процесів у живих організмах (рис.~\ref{pic:interactions}).

\begin{figure}[h!]\centering
	\includegraphics[width=\linewidth]{forces}
	\caption{Взаємодії в природі}
	\label{pic:interactions}
\end{figure}


%Основні типи фізичних полів:
%\begin{enumerate}
%	\item Гравітаційне поле забезпечує взаємодію між масами.
%	\item Електромагнітне поле діє між зарядами та струмами.
%	\item Глюонне поле відповідає за сильну взаємодію.
%	\item Поле векторних бозонів опосередковують слабку взаємодію.
%\end{enumerate}

Серед чотирьох фундаментальних взаємодій електромагнітна займає особливе місце завдяки своєму впливу на всі рівні організації матерії --- від атомів до
макроскопічних об’єктів.\footnote{Сьогодні існують припущення щодо існування матерії, яка не бере участі в електромагнітній взаємодії і тому не випромінює та не поглинає світло.
	Така матерія називається темною. Її існування можна виявити через гравітаційний вплив на звичайну матерію.} Вона визначає структуру атомів і молекул, а
також властивості речовини, які ми спостерігаємо у повсякденному житті. Саме
завдяки електромагнітній взаємодії атоми утримуються разом, забезпечуючи існування твердих і рідких тіл.

Електромагнітна взаємодія також є основою для таких явищ, як поширення світла, радіохвиль та інших форм випромінювання. Ці процеси відіграють ключову
роль у передачі енергії та інформації, що є невід’ємною частиною сучасного життя.

Звичайні сили, які ми відчуваємо у повсякденному житті, такі як тиск руки на камінь чи взаємодія підошви із землею, також обумовлені електромагнітною
взаємодією. Коли два тіла стикаються, їхні атоми зближуються на відстані порядку атомних розмірів, і взаємодія між електронами та ядрами цих атомів
визначає характер цього контакту, створюючи силу, яка не дозволяє тілам проникати одне в одне.

Особливістю електромагнітної взаємодії є її величезна сила порівняно з гравітаційною. Навіть невеликий дисбаланс між електричними зарядами може
призвести до ефектів, що перевершують найпотужніші гравітаційні явища. Як яскраво зазначає Річард Фейнман у своїх лекціях:
\begin{tornpage}\small\itshape%\Alexander
	\hspace{\parindent}Якби у вашому тілі або в тілі вашого сусіда (який стоїть від вас на відстані витягнутої руки) електронів виявилося б всього на 1\%
	більше, ніж протонів, то сила вашого відштовхування була б неймовірно великою. Наскільки великою? Достатньою, щоб підняти хмарочос? Більше!
	Достатньою,
	щоб підняти гору Еверест? Більше! Сили відштовхування вистачило б, щоб підняти «вагу», що дорівнює вазі нашої Землі!\\

	\hfill{ Р. Фейнман у книзі \enquote{Фейнманівські лекції з фізики}}.
\end{tornpage}

%%% --------------------------------------------------------
\section{Електричний заряд}\hypertarget{electric-charge}{}
%%% --------------------------------------------------------


\emphx{Електричний заряд} --- це фізична величина, яка визначає як силу, так і характер (притягання чи відштовхування) взаємодії між зарядженими тілами.
Це базова властивість елементарних частинок, яка не залежить від інших їхніх властивостей і не може бути виведена з якихось глибших принципів. Тому
властивості заряду ми змушені встановлювати не теоретично, а експериментально. Зокрема, дослідним шляхом було виявлено такі його властивості:

\begin{enumerate}
	\item \emphx{Дискретність заряду}.

	      Експеримент, проведений у 1909 році Робертом Ендрюсом Міллікеном та Гарві Флетчером, дозволив встановити, що заряд має дискретну природу. В
	      результаті досліджень було визначено величину заряду електрона з точністю до 1\%, що підтвердило існування мінімальної одиничної величини заряду.

	      Згідно з цими результатами, будь-який ненульовий заряд є кратним деякому елементарному заряду, числове значення якого дорівнює заряду електрона:

	      \begin{equation*} e = 4{,}803 \cdot 10^{-10}\ \text{Фр} = 1{,}602 \cdot 10^{-19}\ \text{Кл}. \end{equation*}

	      Це означає, що заряд не може набувати довільних значень, а лише кратних елементарному заряду.

	      \begin{Historical}
		      Крім того, було встановлено, що в природі існують \emphx{заряди двох типів}. Два роди електрики («скляна» і «смоляна») виявив ще Ш.~Дюфе
		      (1733~р.), а Б.~Франклін запропонував називати їх \emphx{позитивним} та \emphx{негативним} --- ця термінологія й закріпилась у фізиці.
		      Позитивний заряд несуть, наприклад, протони, а негативний --- електрони.
	      \end{Historical}

	      %    \item Адитивність заряду. Заряд є адитивною величиною. Це означає

	\item \emphx{Збереження заряду}.

	      Заряд --- це величина, що \emphx{зберігається}. Позитивні й негативні заряди можуть народжуватися в рівних кількостях або взаємно
	      знищувати один одного. При цьому сумарний електричний заряд замкненої системи не змінюється.
	\item \emphx{Релятивістська інваріантність заряду}.

	      Величина електричного заряду однакова відносно різних систем відліку. Іншими словами, величина заряду не залежить від швидкості.
\end{enumerate}


%% --------------------------------------------------------
\section{Як заряджають тіла?}\hypertarget{charging-bodies}{}
%% --------------------------------------------------------


Тіла складаються з заряджених частинок, і в нормальних умовах їхня кількість є збалансованою, що робить тіло електрично нейтральним. Однак, коли два
тіла вступають у контакт, між ними може відбутися перенесення електронів, що змінює баланс зарядів. Цей процес інтенсифікується при терті між тілами, що
призводить до ще більшого перенесення електронів. Таке явище називається \emphx{трибоелектричним ефектом} (рис.~\ref{pic:triboeffect}). В результаті
цього ефекту одне тіло набуває негативного заряду, а інше --- позитивного.

%---------------------------------------------------------
\begin{SCfigure}[][h!]
	\centering
	\includegraphics[height=3cm]{charging-by-friction}
	\caption{Демонстрація трибоелектричного ефекту.\label{pic:triboeffect}}
\end{SCfigure}
%---------------------------------------------------------

Слово \enquote{\emphx{електрика}} походить від грецького \enquote{$\eta\lambda\varepsilon\kappa\tau\rho o\nu$} (\emphx{електрон}), що означає \enquote{\emphx{бурштин}}. Люди помітили, що бурштин, натертий об хутро, притягує дрібні предмети, і це явище згодом стало основою для дослідження електричних явищ.

Матеріали, що виявляють трибоелектричний ефект, прийнято розташовувати в \emphx{трибоелектричний ряд}, один кінець якого є позитивним, а інший --- негативним. Під час тертя пари матеріалів з ряду, матеріал, розташований ближче до позитивного кінця ряду, зарядиться позитивно, а інший --- негативно (рис.~\ref{tikz:triboseries}).


\begin{figure}[h!]
	\centering
	\begin{tikzpicture}
		% Градієнт фону
		\shade[left color=blue!30, right color=red!30] (-5.5, -.5) rectangle (5.5, -3.5);

		\node[single  arrow, single  arrow head extend=.5cm, left color=blue!50, right color=red!50, text=white, font={\bf \large}, inner sep=0.3cm]
		(doublearrow)
		{--\hspace*{10.5cm}+};

		% Трибоелектричний ряд
		\node[align=center, rotate=90] at (5,  -2) {Хутро};
		\node[align=center, rotate=90] at (4,  -2) {Фланель};
		\node[align=center, rotate=90, align=center, text width=2.5cm] at (3,  -2) {Слонова\\ кістка};
		\node[align=center, rotate=90] at (2,  -2) {Пір'я};
		\node[align=center, rotate=90] at (1,  -2) {Флінтглас};
		\node[align=center, rotate=90] at (0,  -2) {Папір};
		\node[align=center, rotate=90] at (-1, -2) {Шовк};
		\node[align=center, rotate=90] at (-2, -2) {Алюміній};
		\node[align=center, rotate=90] at (-3, -2) {Бурштин};
		\node[align=center, rotate=90] at (-4, -2) {Ебоніт};
		\node[align=center, rotate=90] at (-5, -2) {Поліестер};
	\end{tikzpicture}
	\caption{Трибоелектричний ряд\label{tikz:triboseries}}
\end{figure}


Трибоелектричний ефект має не лише пізнавальне, а й практичне значення. З одного боку, він лежить в основі \emphx{трибоелектричних наногенераторів} (TENG) --- пристроїв, які перетворюють механічну енергію тертя чи вібрацій (кроки людини, вітер, хвилі) в електричну; завдяки високій ефективності при низьких частотах їх розглядають як перспективне джерело живлення для носимої електроніки та автономних сенсорів у системах \enquote{інтернету речей}.

З іншого боку, неконтрольоване трибоелектричне заряджання створює практичні проблеми. У мікроелектроніці раптовий електростатичний розряд (\emphx{ESD}) здатний вивести з ладу чутливі напівпровідникові елементи, тому виробництво мікросхем супроводжується суворими заходами захисту від статичної електрики (антистатичні браслети, заземлені робочі поверхні, іонізатори повітря). У хімічній та харчовій промисловості трибоелектричне заряджання дрібнодисперсних матеріалів (борошна, вугільного пилу) може спричинити електростатичний розряд --- одну з причин вибухів пилоповітряних сумішей. Водночас керований трибоелектричний ефект використовують і на користь: в \emphx{електрофільтрах} заряджені частинки диму чи пилу осідають на протилежно заряджених електродах, очищуючи промислові гази, а в \emphx{ксерографії} (принцип роботи копіювальних апаратів і лазерних принтерів) електростатичне заряджання світлочутливого барабана й тонера забезпечує перенесення зображення на папір.


%% --------------------------------------------------------
\section{Закон Кулона}\hypertarget{coulombs-law}{}
%% --------------------------------------------------------


Експериментально встановлено, що заряди одного знака (однойменні заряди) \emphx{відштовхуються}, а заряди різного знака (різнойменні заряди) --- \emphx{притягуються}. Ці властивості зарядів лягли в основу формулювання закону їхньої взаємодії.


Для спрощення аналізу часто використовують ідеалізовані моделі заряджених тіл, однією з яких є \emphx{точковий заряд}. Точковим називається такий заряд, розмірами якого в конкретних умовах можна знехтувати.

Із цією моделлю заряджених тіл пов'язаний важливий закон, встановлений Ш.~Кулоном у 1785 році. Він стверджує, що \emphx{сила, що діє між двома нерухомими точковими зарядами $q_1$ і $q_2$, розташованими на відстані $r$ один від одного, описується законом}:
\begin{equation}\label{eq:ColumbLaw}
	\tcbhighmath{
		\vect{F}_{12} = \frac{q_1 q_2}{r^2}\frac{\vect{r}_{12}}{r}.
	}
\end{equation}
Ця сила напрямлена вздовж лінії, яка з'єднує заряди, тобто носить \emphx{центральний} характер.
\begin{Attention}
	У реальних умовах точкові заряди не існують, і для застосування закону Кулона до малих заряджених тіл необхідно, щоб заряд всередині них був розподілений рівномірно, інакше порушується центральний характер взаємодії. Тому під \enquote{фізично} точковими зарядами будемо розуміти такі, в яких розподіл заряду є рівномірним, тобто не залежить від орієнтації тіла.
\end{Attention}

На рис.~\ref{pic:Columb} показані величини, що входять до закону Кулона, підкреслено центральний характер взаємодії та її напрям однойменних та різнойменних зарядів.


\begin{SCfigure}[0.5][h!]
	\centering
	\localinput{Columb.tikz}
	\caption{Ілюстрація до закону Кулона}
	\label{pic:Columb}
\end{SCfigure}

Досліди самого Кулона не відзначалися високою точністю, і можна було припустити, що знайдений ним простий закон обернених квадратів виконується лише наближено, наприклад:
\begin{equation*}
	F = \frac{q_1 q_2}{r^{2 + \epsilon}},
\end{equation*}
де $|\epsilon| \ll 1$.

У випадку, коли $\epsilon = 0$, заряди зосереджуються лише на поверхні провідника, а всередині провідника їх не існує. Якщо припустити $|\epsilon| \neq 0$, то це може спричинити появу заряду в товщі провідника.\footnote{Строге обґрунтування цього факту спирається на теорему Гаусса, з якою ми познайомимось у розділі~\ref{SteadyEfield}, \enquote{\nameref{sec:Gauss}}; поки що приймемо його як наслідок точного закону обернених квадратів.}

\begin{SCfigure}[0.75][h!]
	\centering
	\includegraphics[width=0.5\linewidth]{cavindesh_exp}
	\caption{Установка Г.~Кавендіша для перевірки закону обернених квадратів.}
	\label{pic:Cavendih_exp}
\end{SCfigure}

Експеримент, що дозволив перевірити ці припущення, був проведений Г.~Кавендішем і полягав у накладанні двох провідних напівсфер на заряджену провідну кулю, що щільно прилягали одна до одної, утворюючи майже суцільну оболонку (рис.~\ref{pic:Cavendih_exp}). За умов закону обернених квадратів весь заряд мав би перейти на зовнішні напівсфери. Якщо ж взаємодія описується за формулою з $\epsilon \neq 0$, то всередині кулі залишався б ненульовий заряд, який пропорційно зменшується зі зменшенням значення $\epsilon$. Цей заряд можна було б виміряти за допомогою чутливого електрометра, приєднаного до внутрішньої сфери (на рисунку роль електрометра відіграють дві кульки на нитках). Експеримент повторювався з поступовим підвищенням точності вимірювань. У 1983 році було отримано обмеження для $\epsilon$:
\begin{equation*}
	|\epsilon| < 10^{-16} \text{–} 10^{-17}.
\end{equation*}

\begin{Attention}
	Порушення закону обернених квадратів мало б серйозні наслідки для теорії електромагнітної взаємодії, оскільки цей закон є основою для її опису. Теорія Максвелла, яка визначає поведінку електричних та магнітних полів, спирається на його точність. Якби закон порушувався, це змінило б наші уявлення про електромагнітну взаємодію та вимагало б перегляду численних експериментальних результатів, що сьогодні підтверджують цю теорію.
\end{Attention}


%% --------------------------------------------------------
\section*{Підсумки}
\phantomsection
\addcontentsline{toc}{section}{Підсумки}
%% --------------------------------------------------------

Серед чотирьох фундаментальних взаємодій саме електромагнітна визначає будову атомів і молекул, властивості речовини та всі \enquote{звичайні} сили повсякденного життя. Її джерелом є електричний заряд --- фундаментальна величина, властивості якої встановлено з досліду: заряд \emphx{дискретний} (кратний елементарному заряду $e = 4{,}803\cdot10^{-10}$~Фр), \emphx{зберігається} в замкненій системі та \emphx{релятивістськи інваріантний}. Тіла заряджаються перенесенням електронів, зокрема при терті (трибоелектричний ефект), а електрична нейтральність макроскопічних тіл є наслідком майже точної рівноваги позитивних і негативних зарядів.

Кількісно взаємодію двох нерухомих точкових зарядів описує закон Кулона~\eqref{eq:ColumbLaw}: сила є центральною, обернено пропорційною квадрату відстані, і відштовхує однойменні заряди та притягує різнойменні. Точність закону обернених квадратів підтверджена дослідами типу кавендішевого: $|\epsilon| < 10^{-16}\text{--}10^{-17}$. Закон Кулона, однак, лише фіксує результат взаємодії двох зарядів, нічого не кажучи про те, \emphx{як} вона передається. Наступний логічний крок --- ввести поняття \emphx{електричного поля} як посередника взаємодії, навчитися знаходити поле системи зарядів і, як наслідок, довести теорему Гаусса (гл.~\ref{SteadyEfield}).

%% --------------------------------------------------------
\section*{Задачі}
\phantomsection
\addcontentsline{toc}{section}{Задачі}
%% --------------------------------------------------------


\begin{problem}\label{prob:FeynmanEstimate}
	\emphx{Перевіряємо Фейнмана.} Оцініть силу відштовхування двох людей масою $70$~кг кожна на відстані $1$~м, якщо в кожного електронів на $1\%$ більше, ніж протонів. Порівняйте її з вагою Землі ($M_\oplus\approx6\cdot10^{27}$~г). Вважайте, що приблизно половина нуклонів тіла --- протони, $m_p\approx1{,}7\cdot10^{-24}$~г.
	%	\emphx{Відповідь:} $N_p\approx\dfrac{7\cdot10^{4}}{2\cdot1{,}7\cdot10^{-24}}\approx2\cdot10^{28}$, надлишковий заряд $q\approx2\cdot10^{26}\cdot4{,}8\cdot10^{-10}\approx10^{17}$~Фр, $F=q^2/r^2\approx10^{30}$~дин $\approx10^{25}$~Н. Вага Землі $M_\oplus g\approx6\cdot10^{30}$~дин, тобто сила того ж порядку.
\end{problem}

\begin{problem}\label{prob:ChargeSharing}
	\emphx{Дотик кульок.} Дві однакові провідні кульки із зарядами $+3q$ і $-q$ розташовані на відстані $r\gg R$ одна від одної. Їх привели в дотик і розвели на ту саму відстань. У скільки разів і як змінилася сила взаємодії?
	%	\emphx{Відповідь:} За збереженням заряду та симетрією на кожній кульці стало по $+q$. Було притягання $F_1=3q^2/r^2$, стало відштовхування $F_2=q^2/r^2$: сила зменшилась у $3$ рази й змінила напрям.
\end{problem}

\begin{problem}\label{prob:CoulombVsGravity}
	\emphx{Електрика проти гравітації.} Знайдіть відношення кулонівської сили до гравітаційної для електрона й протона. Чи залежить воно від відстані між ними? ($e=4{,}8\cdot10^{-10}$~Фр, $m_e=9{,}1\cdot10^{-28}$~г, $m_p=1{,}67\cdot10^{-24}$~г, $G=6{,}67\cdot10^{-8}\ \text{см}^3/(\text{г}\cdot\text{с}^2)$.)
	\emphx{Відповідь:} $\dfrac{F_e}{F_g}=\dfrac{e^2}{Gm_em_p}\approx2\cdot10^{39}$; від відстані не залежить, бо обидві сили $\propto1/r^2$.
\end{problem}


%% --------------------------------------------------------
\begin{Questions}
	\begin{quizlist}
		\item Назвіть чотири фундаментальні взаємодії. Яка з них визначає будову атомів і молекул та силу, з якою рука тисне на камінь? Чому цю силу вважають електромагнітною за природою?
		\item Чому в повсякденному житті ми не відчуваємо величезних електричних сил, про які говорить Фейнман? Що забезпечує електричну нейтральність макроскопічних тіл?
		\item Що означає, що електричний заряд є фундаментальною величиною? Чому його властивості ми вивчаємо з досліду, а не виводимо з інших принципів?
		\item Сформулюйте три основні властивості заряду: дискретність, збереження та релятивістську інваріантність. Чи може тіло мати заряд $2{,}5\,e$? Відповідь обґрунтуйте.
		\item У чому полягав дослід Міллікена та Флетчера і що він довів? Який зв'язок між значеннями елементарного заряду в одиницях СГС (Фр) та СІ (Кл)?
		\item Що таке трибоелектричний ефект і чому при терті двох тіл одне заряджається позитивно, а інше --- негативно? Користуючись трибоелектричним рядом (рис.~\ref{tikz:triboseries}), визначте знаки зарядів, яких набудуть шовк і поліестер при терті один об одного.
		\item Наведіть по одному прикладу корисного та шкідливого прояву трибоелектричного ефекту. Поясніть, як працює електрофільтр або ксерографія.
		\item Сформулюйте закон Кулона й запишіть його у векторній формі. Що означає, що кулонівська сила є \emphx{центральною}?
		\item За яких умов заряджене тіло можна вважати точковим? Чому для «фізично» точкових зарядів потрібен рівномірний розподіл заряду?
		\item Як зміниться сила взаємодії двох точкових зарядів, якщо: а)~відстань між ними збільшити втричі; б)~один із зарядів збільшити вдвічі; в)~обидва заряди змінити знак на протилежний?
		\item У чому полягала ідея досліду Кавендіша? Чому за закону обернених квадратів ($\epsilon = 0$) заряд зосереджується лише на зовнішній поверхні провідника, а при $\epsilon \neq 0$ з'явився б заряд усередині?
		\item Яке обмеження на $|\epsilon|$ отримано в експериментах і чому порушення закону обернених квадратів було б настільки важливим для електродинаміки?
	\end{quizlist}
\end{Questions}